% Results and Discussion (merged §4). Supersedes results002.tex and discuss002.tex.
% \section and \label live in main.tex, per this repo's convention for section files.

Monthly compaction estimates were evaluated under three monitoring conditions: temporary delay of complete monthly MLCW records, periodic reduction of MLCW measurement frequency to once every six or twelve months, and permanent cessation of MLCW reporting after 3, 5, or 8 years. Each condition removes a different amount of direct subsurface information, and the three conditions are reported separately because they draw on different frozen model snapshots and answer different operational questions.

\subsection{Overall nowcasting performance and depth-dependence}
\label{subsec:results_overall_performance}

The delayed-data evaluation produced 138 monthly estimates per depth section across 23 complete six-month blocks, for a pooled total of 828 estimates across the six sections. Performance varied systematically with depth (\Cref{tab:delayed_performance}). Pooled across all six sections, the model reached $R^2=0.778$, RMSE~$=0.423$~mm/month, and MAE~$=0.279$~mm/month. Section-level $R^2$ ranged from 0.076 at the 200--250~m section (S5) to 0.976 at the second section (S2), with S1, S3, and S4 falling between these two ($R^2=0.888$, 0.870, and 0.798, respectively) and S6 also weak ($R^2=0.337$).

The weak result at S5 traces to a physical observability gap rather than a modelling failure. The groundwater-level record used as a predictor for S5 is drawn from a well screened at 257--263~m, entirely below the 200--250~m depth range that S5 represents, and the same well also supplies the predictor for the adjacent S6 section. No piezometer is screened within S5's own compacting depth interval, so the head-change predictor available for that section is an imprecise proxy for the pore-pressure conditions actually driving compaction there. If a station's monitoring network is reduced by removing sensors from an actively compacting layer, the ability to nowcast that layer is lost regardless of how much surface or head data remain elsewhere in the profile.

Each estimate also carried a 90\% Bayesian posterior predictive interval, available from the first evaluation block onward because the interval is calculated directly from the fitted posterior rather than from an archive of past errors. Pooled across all sections, 78.0\% of observations fell within their nominal 90\% interval, a shortfall of about twelve percentage points below the nominal level; the interval was therefore narrower than the data actually required. Coverage varied by section, from 65.9\% at S6 to 88.4\% at S2, and did not track $R^2$ directly: S5, the weakest section by point accuracy, reached 81.2\% coverage with a mean interval width of 1.85~mm/month, more than twice the width at any other section, so its comparatively wide interval partly compensated for its weak point estimate. \Cref{subsec:results_limitations} returns to the interpretation of this shortfall.

\begin{table}[htbp]
  \centering
  \scriptsize
  \renewcommand{\arraystretch}{1.12}
  \caption{Performance and posterior predictive interval statistics of monthly compaction estimates during temporary MLCW data delays (mm/month unless noted).}
  \label{tab:delayed_performance}
  \resizebox{\textwidth}{!}{%
  \begin{tabular}{lrrrrrr}
    \toprule
    \textbf{Section} & \textbf{$n$} & \textbf{$R^2$} & \textbf{RMSE} & \textbf{MAE} & \textbf{Coverage (\%)} & \textbf{Width} \\
    \midrule
    S1 & 138 & 0.888 & 0.212 & 0.178 & 74.6 & 0.541 \\
    S2 & 138 & 0.976 & 0.239 & 0.182 & 88.4 & 0.637 \\
    S3 & 138 & 0.870 & 0.229 & 0.153 & 73.9 & 0.443 \\
    S4 & 138 & 0.798 & 0.311 & 0.259 & 84.1 & 0.985 \\
    S5 & 138 & 0.076 & 0.661 & 0.504 & 81.2 & 1.848 \\
    S6 & 138 & 0.337 & 0.621 & 0.397 & 65.9 & 0.972 \\
    \midrule
    All & 828 & 0.778 & 0.423 & 0.279 & 78.0 & 0.904 \\
    \bottomrule
  \end{tabular}%
  }
\end{table}

Because each estimate used only hydraulic head and cGNSS observations available in the same target month, this capability is best understood as same-month nowcasting rather than forecasting. The model's primary operational value is bridging the interval between the target month and the delayed delivery of the corresponding MLCW measurement, so that pumping-management decisions do not have to wait for the historical record to be reconciled.

\subsection{Coefficients of driving factors at each section}
\label{subsec:results_driving_factors}

The predictors that dominated the fit differed by depth section, identified as the single feature with the largest-magnitude median standardized coefficient across folds at each section (\Cref{tab:driving_factors}). A seasonal term was the dominant predictor at five of the six sections: \texttt{month\_sin} at S1 (median $-0.212$) and S2 (median $-0.542$), \texttt{month2\_cos} at S3 (median $-1.256$), \texttt{month\_sin} at S5 (median $-0.072$), and \texttt{month\_cos} at S6 (median $+0.310$). The single exception was S4, where the dominant predictor was the cGNSS surface displacement term (median $+0.236$), not a seasonal or hydraulic-head term.

This single-feature ranking rule structurally favours seasonal terms over hydraulic-head terms because of how each predictor group is constructed: only five seasonal features were available to any section, against 25 hydraulic-head features per section (ten describing the target section's own screened interval, fifteen describing the other five monitored depth intervals as candidate system-wide predictors). A single dominant seasonal feature can therefore carry more concentrated weight than a hydraulic-head signal that is instead spread across many correlated head-change terms. The five-of-six seasonal dominance pattern reported here reflects this rule and would not necessarily hold under a rule that summed coefficient magnitude within each predictor group rather than selecting the single largest feature.

The cross-section hydraulic-head terms, included as candidate predictors representing system-wide hydraulic conditions at Tuku, did not dominate at any of the six sections, indicating that each section's own local head record, when it existed, carried more concentrated predictive weight than head changes observed at the other five depth intervals.

\begin{table}[htbp]
  \centering
  \small
  \renewcommand{\arraystretch}{1.15}
  \caption{Dominant predictor group at each depth section, identified as the single feature with the largest-magnitude median standardized coefficient across folds.}
  \label{tab:driving_factors}
  \begin{tabular}{lllr}
    \toprule
    \textbf{Section} & \textbf{Dominant group} & \textbf{Feature} & \textbf{Median coefficient} \\
    \midrule
    S1 & Seasonality & \texttt{month\_sin} & $-0.212$ \\
    S2 & Seasonality & \texttt{month\_sin} & $-0.542$ \\
    S3 & Seasonality & \texttt{month2\_cos} & $-1.256$ \\
    S4 & Deformation (cGNSS) & \texttt{dS\_total} & $+0.236$ \\
    S5 & Seasonality & \texttt{month\_sin} & $-0.072$ \\
    S6 & Seasonality & \texttt{month\_cos} & $+0.310$ \\
    \bottomrule
  \end{tabular}
\end{table}

\subsection{Sensitivity to reduced-frequency measurements}
\label{subsec:results_reduced_frequency}

Reducing the frequency of MLCW field visits from monthly to once every six or twelve months increased endpoint error at every section and every initial calibration length tested (\Cref{tab:reduced_frequency}). Pooled across all sections, endpoint MAE rose from approximately 1.4--1.5~mm under the six-month schedule to approximately 2.4--2.7~mm under the twelve-month schedule, roughly a 70--80\% increase, depending on the initial calibration length. The increase was not uniform across sections: at S5, endpoint MAE rose from 2.65~mm (six-month schedule, 36-month initial history) to 4.75~mm (twelve-month schedule, same initial history), while at S1 the corresponding increase was from 0.94~mm to 1.06~mm, a much smaller absolute change. For the thinnest schedule evaluated, a twelve-month interval following 96 months of initial monthly calibration, six complete endpoints were available per section within the common record length.

\begin{table}[htbp]
  \centering
  \scriptsize
  \renewcommand{\arraystretch}{1.12}
  \caption{Endpoint mean absolute error (mm) under six- and twelve-month reduced MLCW measurement schedules, by depth section and initial calibration length.}
  \label{tab:reduced_frequency}
  \begin{tabular}{lrrr}
    \toprule
    \textbf{Section} & \textbf{Initial history (months)} & \textbf{MAE, 6-month schedule} & \textbf{MAE, 12-month schedule} \\
    \midrule
    S1 & 36 & 0.94 & 1.06 \\
    S1 & 96 & 1.02 & 1.33 \\
    S5 & 36 & 2.65 & 4.75 \\
    S5 & 96 & 2.55 & 4.76 \\
    \midrule
    All & 36 & 1.47 & 2.46 \\
    All & 96 & 1.52 & 2.74 \\
    \bottomrule
  \end{tabular}
\end{table}

This trade-off frames a direct operational choice for a monitoring agency. Halving the field-visit frequency, from every six months to every twelve, roughly doubles the endpoint error at the sections that already estimate compaction least reliably, while sections with a well-constrained hydraulic-head record absorb the same reduction with a comparatively small increase in error. A schedule decision made without regard to depth section would apply the same visit frequency to sections with very different sensitivity to that reduction.

\subsection{Sensitivity to permanent monitoring stoppage}
\label{subsec:results_permanent_stoppage}

Fitting the model once on 3, 5, or 8 years of initial calibration data and then estimating continuously thereafter with no refit and no further MLCW input produced cumulative absolute error that grew with an oscillating, not strictly monotonic, trajectory: for the 3-year training window at S1, 81 of 139 month-to-month changes in cumulative absolute error were increases and 58 were decreases, consistent with partial cancellation of signed monthly errors superimposed on a slowly rising envelope rather than steady monotonic drift.

At a common horizon of 80 months after the training window ended, the shared endpoint across all three training-window lengths, cumulative absolute error averaged across the six sections was 17.0~mm for the 3-year window, 8.3~mm for the 5-year window, and 4.3~mm for the 8-year window (\Cref{tab:permanent_stoppage}). The worst individual cell was the 3-year window at S5 (53.5~mm at the 80-month horizon), consistent with the same observability gap identified in \Cref{subsec:results_overall_performance}: the well supplying S5's head predictor is screened below S5's own depth interval, and a model with less initial calibration data to constrain that already-imprecise relation drifted furthest once no further MLCW correction was possible. The best individual cell was the 8-year window at S3 (0.1~mm at the 80-month horizon).

These three training-window lengths cover different calendar periods of the monitoring record as well as different amounts of initial data: the 3-year window began estimating earlier in the record than the 8-year window, so the improvement from 3 to 8 years reflects both the additional calibration data and the different hydrogeological conditions encountered during each window's subsequent estimation period. The pattern is reported here as three distinct operational scenarios, not as isolated evidence that a longer training window causes better long-term extrapolation.

\begin{table}[htbp]
  \centering
  \small
  \renewcommand{\arraystretch}{1.15}
  \caption{Cumulative absolute error (mm) at a common 80-month post-training-window horizon, averaged across the six depth sections, for three permanent-stoppage training-window lengths.}
  \label{tab:permanent_stoppage}
  \begin{tabular}{lr}
    \toprule
    \textbf{Training window} & \textbf{Mean cumulative absolute error at 80 months (mm)} \\
    \midrule
    3 years & 17.0 \\
    5 years & 8.3 \\
    8 years & 4.3 \\
    \bottomrule
  \end{tabular}
\end{table}

\subsection{Limitations and practical scope}
\label{subsec:results_limitations}

The fitted model is strictly local to Tuku's monitoring configuration and stress history recorded during the evaluation period. The methodology, comprising the per-section walk-forward regression, the posterior predictive interval, and the reduced-frequency and permanent-stoppage sensitivity designs, is transferable to another monitoring station, but the specific fitted coefficients reported here are not, because they depend on the particular hydraulic-head wells and cGNSS record available at Tuku.

The posterior predictive interval reported throughout \Cref{subsec:results_overall_performance,subsec:results_permanent_stoppage} is a Bayesian posterior interval calculated directly from the fitted model, not a split-conformal interval calibrated from an archive of past prediction errors. Split-conformal calibration was implemented for a related, pooled formulation of the model but was not executed for the depth-section-independent formulation reported here. The pooled shortfall below nominal coverage reported in \Cref{subsec:results_overall_performance} therefore describes this posterior interval alone, and is not evidence about how a genuinely calibrated conformal interval would have performed had one been computed for this formulation.

The reduced-frequency and permanent-stoppage sensitivity analyses in \Cref{subsec:results_reduced_frequency,subsec:results_permanent_stoppage} share one snapshot of the fitted model and are directly comparable with each other, but neither is directly comparable, number for number, with the delayed-data results in \Cref{subsec:results_overall_performance}, because the two groups of analyses were built and frozen from different underlying model snapshots. Any comparison across this boundary should be stated as a qualitative contrast between monitoring conditions, not as a numeric difference between two runs of the same fitted model.

Surface displacement from cGNSS integrates deformation above and below the deepest MLCW ring rather than resolving it by depth, and each groundwater well observes hydraulic head only within its own screened interval. The lithological log provided geological context for comparing the depth sections but did not directly measure material properties such as compressibility or preconsolidation stress. Physical attribution of the mechanism behind any single month's compaction, beyond the statistical association reported here, would require geotechnical measurements or a coupled groundwater flow and aquifer-system compaction model \citep{galloway_land_1999, hung2012_mlcw}.

\FloatBarrier
